\begin{thebibliography}{99}
\setlength{\itemsep}{2pt plus 0.5pt}\setlength{\parsep}{0pt}
\bibitem{ref1}\label{r1} Callen JL, Westbrook JI, Georgiou A, et al. Failure to follow-up test results for ambulatory patients: a systematic review. J Gen Intern Med 2012;27(10):1334--1348. doi:10.1007/s11606-011-1949-5.
\bibitem{ref2}\label{r2} Patel MP, Schettini P, O'Leary CP, et al. Closing the Referral Loop: an Analysis of Primary Care Referrals to Specialists in a Large Health System. J Gen Intern Med 2018;33(5):715--721. doi:10.1007/s11606-018-4392-z.
\bibitem{ref3}\label{r3} Kwan JL, Singh H. Assigning responsibility to close the loop on radiology test results. Diagnosis (Berl) 2017;4(3):173--177. doi:10.1515/dx-2017-0019.
\bibitem{ref4}\label{r4} Singh H, Meyer AND, Thomas EJ. The frequency of diagnostic errors in outpatient care. BMJ Qual Saf 2014;23(9):727--731. doi:10.1136/bmjqs-2013-002627.
\bibitem{ref5}\label{r5} Slawomirski L, Kelly D, de Bienassis K, et al. The economics of diagnostic safety. OECD Health Working Papers 2025. doi:10.1787/fc61057a-en.
\bibitem{ref6}\label{r6} Wright B, Lennox A, Graber ML, et al. Closing the loop on test results to reduce communication failures: a rapid review of evidence, practice and patient perspectives. BMC Health Serv Res 2020;20:897. doi:10.1186/s12913-020-05737-x.
\bibitem{ref7}\label{r7} Duan-Porter W, Ullman K, Majeski B, et al. Care Coordination Models and Tools: Systematic Review and Key Informant Interviews. J Gen Intern Med 2022;37(6):1367--1379. doi:10.1007/s11606-021-07158-w.
\bibitem{ref8}\label{r8} Shen Y, Yu J, Zhou J, et al. Twenty-Five Years of Evolution and Hurdles in Electronic Health Records and Interoperability in Medical Research: Comprehensive Review. J Med Internet Res 2025;27:e59024. doi:10.2196/59024.
\bibitem{ref9}\label{r9} Evans RS. Electronic Health Records: Then, Now, and in the Future. Yearb Med Inform 2016;25(Suppl 1):S48--S61. doi:10.15265/IYS-2016-s006.
\bibitem{ref10}\label{r10} Rowley J. The wisdom hierarchy: representations of the DIKW hierarchy. J Inf Sci 2007;33(2):163--180. doi:10.1177/0165551506070706.
\bibitem{ref11}\label{r11} Li R, Niu Y, Scott SR, et al. Using Electronic Medical Record Data for Research in a HIMSS Analytics EMRAM Stage 7 Hospital in Beijing: Cross-sectional Study. JMIR Med Inform 2021;9(8):e24405. doi:10.2196/24405.
\bibitem{ref12}\label{r12} Duncan R, Eden R, Woods L, et al. Synthesizing Dimensions of Digital Maturity in Hospitals: Systematic Review. J Med Internet Res 2022;24(3):e32994. doi:10.2196/32994.
\bibitem{ref13}\label{r13} Blumenthal D, Tavenner M. The ``Meaningful Use'' Regulation for Electronic Health Records. N Engl J Med 2010;363:501--504. doi:10.1056/NEJMp1006114.
\bibitem{ref14}\label{r14} Adler-Milstein J, DesRoches CM, Kralovec P, et al. Electronic Health Record Adoption In US Hospitals: Progress Continues, But Challenges Persist. Health Aff 2015;34(12):2174--2180. doi:10.1377/hlthaff.2015.0992.
\bibitem{ref15}\label{r15} Weed LL. Medical Records That Guide and Teach. N Engl J Med 1968;278:593--600. doi:10.1056/NEJM196803212781204.
\bibitem{ref16}\label{r16} Bodenreider O, Cornet R, Vreeman DJ. Recent Developments in Clinical Terminologies: SNOMED CT, LOINC, and RxNorm. Yearb Med Inform 2018;27(1):129--139. doi:10.1055/s-0038-1667077.
\bibitem{ref17}\label{r17} Bender D, Sartipi K. HL7 FHIR: An Agile and RESTful Approach to Healthcare Information Exchange. IEEE CBMS 2013:326--331. doi:10.1109/CBMS.2013.6627810.
\bibitem{ref18}\label{r18} Mandel JC, Kreda DA, Mandl KD, et al. SMART on FHIR: a standards-based, interoperable apps platform for electronic health records. J Am Med Inform Assoc 2016;23(5):899--908. doi:10.1093/jamia/ocv189.
\bibitem{ref19}\label{r19} Friedman C, Alderson PO, Austin JHM, et al. A general natural-language text processor for clinical radiology (MedLEE). J Am Med Inform Assoc 1994;1(2):161--174. doi:10.1136/jamia.1994.95236146.
\bibitem{ref20}\label{r20} Savova GK, Masanz JJ, Ogren PV, et al. Mayo clinical Text Analysis and Knowledge Extraction System (cTAKES). J Am Med Inform Assoc 2010;17(5):507--513. doi:10.1136/jamia.2009.001560.
\bibitem{ref21}\label{r21} Safran C, Bloomrosen M, Hammond WE, et al. Toward a National Framework for the Secondary Use of Health Data: An AMIA White Paper. J Am Med Inform Assoc 2007;14(1):1--9. doi:10.1197/jamia.M2273.
\bibitem{ref22}\label{r22} Rosenbloom ST, Denny JC, Xu H, et al. Data from clinical notes: a perspective on the tension between structure and flexible documentation. J Am Med Inform Assoc 2011;18(2):181--186. doi:10.1136/jamia.2010.007237.
\bibitem{ref23}\label{r23} HL7 FHIR Release 5: Workflow. Health Level Seven International; 2023. https://hl7.org/fhir/R5/workflow.html (accessed 4 Sep 2026).
\bibitem{ref24}\label{r24} Peleg M. Computer-interpretable clinical guidelines: a methodological review. J Biomed Inform 2013;46(4):744--763. doi:10.1016/j.jbi.2013.06.009.
\bibitem{ref25}\label{r25} Lenz R, Reichert M. IT support for healthcare processes -- premises, challenges, perspectives. Data Knowl Eng 2007;61(1):39--58. doi:10.1016/j.datak.2006.04.007.
\bibitem{ref26}\label{r26} Aspland E, Gartner D, Harper P. Clinical pathway modelling: a literature review. Health Syst 2021;10(1):1--23. doi:10.1080/20476965.2019.1652547.
\bibitem{ref27}\label{r27} van der Aalst W. Process Mining: Data Science in Action. 2nd ed. Berlin: Springer; 2016. doi:10.1007/978-3-662-49851-4.
\bibitem{ref28}\label{r28} Rojas E, Munoz-Gama J, Sepúlveda M, et al. Process mining in healthcare: A literature review. J Biomed Inform 2016;61:224--236. doi:10.1016/j.jbi.2016.04.007.
\bibitem{ref29}\label{r29} Office of the National Coordinator for Health Information Technology (ONC), HHS. 21st Century Cures Act: Interoperability, Information Blocking, and the ONC Health IT Certification Program. Final rule. Fed Regist 2020;85(85):25642--25961. federalregister.gov/d/2020-07419 (accessed 4 Sep 2026). https://www.federalregister.gov/documents/2020/05/01/2020-07419/21st-century-cures-act-interoperability-information-blocking-and-the-onc-health-it-certification
\bibitem{ref30}\label{r30} Griffin AC, He L, Sunjaya AP, et al. Clinical, technical, and implementation characteristics of real-world health applications using FHIR. JAMIA Open 2022;5(4):ooac077. doi:10.1093/jamiaopen/ooac077.
\bibitem{ref31}\label{r31} Corbeil JP, Ben Abacha A, Tremblay J, et al. Overview of the MEDIQA-OE 2025 Shared Task on Medical Order Extraction from Doctor-Patient Consultations. arXiv preprint; 2025. arXiv:2510.26974.
\bibitem{ref32}\label{r32} Kanaparthy NS, Bakare T, Diao Z, et al. Real-World Evidence Synthesis of Digital Scribes Using Ambient Listening and Generative Artificial Intelligence for Clinician Documentation Workflows: Rapid Review. JMIR AI 2025;4:e76743. doi:10.2196/76743.
\bibitem{ref33}\label{r33} Soysal E, Wang J, Jiang M, et al. CLAMP: a toolkit for efficiently building customized clinical NLP pipelines. J Am Med Inform Assoc 2018;25(3):331--336. doi:10.1093/jamia/ocx132.
\bibitem{ref34}\label{r34} Yetisgen-Yildiz M, Gunn ML, Xia F, et al. A text processing pipeline to extract recommendations from radiology reports. J Biomed Inform 2013;46(2):354--362. doi:10.1016/j.jbi.2012.12.005.
\bibitem{ref35}\label{r35} Reichenpfader D, Müller H, Denecke K. A scoping review of large language model based approaches for information extraction from radiology reports. npj Digit Med 2024;7:222. doi:10.1038/s41746-024-01219-0.
\bibitem{ref36}\label{r36} Henry S, Buchan K, Filannino M, et al. 2018 n2c2 shared task on adverse drug events and medication extraction in electronic health records. J Am Med Inform Assoc 2020;27(1):3--12. doi:10.1093/jamia/ocz166.
\bibitem{ref37}\label{r37} Agrawal M, Hegselmann S, Lang H, et al. Large language models are few-shot clinical information extractors. Proc. EMNLP 2022:1998--2022. doi:10.18653/v1/2022.emnlp-main.130.
\bibitem{ref38}\label{r38} Bethard S, Savova G, Palmer M, et al. SemEval-2017 Task 12: Clinical TempEval. Proc. 11th Int Workshop Semantic Eval (SemEval-2017) 2017:565--572. doi:10.18653/v1/S17-2093.
\bibitem{ref39}\label{r39} Mullenbach J, Pruksachatkun Y, Adler S, et al. CLIP: A Dataset for Extracting Action Items for Physicians from Hospital Discharge Notes. Proc. 59th Annu Meet Assoc Comput Linguist (ACL-IJCNLP) 2021:1365--1378. doi:10.18653/v1/2021.acl-long.109.
\bibitem{ref40}\label{r40} Elgaar M, Cheng J, Vakil N, et al. MedDec: A Dataset for Extracting Medical Decisions from Discharge Summaries. Findings Assoc Comput Linguist: ACL 2024 2024:16442--16455. doi:10.18653/v1/2024.findings-acl.975.
\bibitem{ref41}\label{r41} Uzuner O, South BR, Shen S, et al. 2010 i2b2/VA challenge on concepts, assertions, and relations in clinical text. J Am Med Inform Assoc 2011;18(5):552--556. doi:10.1136/amiajnl-2011-000203.
\bibitem{ref42}\label{r42} Sun W, Rumshisky A, Uzuner O. Evaluating temporal relations in clinical text: 2012 i2b2 Challenge. J Am Med Inform Assoc 2013;20(5):806--813. doi:10.1136/amiajnl-2013-001628.
\bibitem{ref43}\label{r43} Strötgen J, Gertz M. Multilingual and cross-domain temporal tagging (HeidelTime). Lang Resour Eval 2013;47(2):269--298. doi:10.1007/s10579-012-9179-y.
\bibitem{ref44}\label{r44} Chang AX, Manning CD. SUTime: A library for recognizing and normalizing time expressions. Proc. LREC 2012:3735--3740.
\bibitem{ref45}\label{r45} Knez T, Žitnik S. Multimodal learning for temporal relation extraction in clinical texts. J Am Med Inform Assoc 2024;31(6):1380--1387. doi:10.1093/jamia/ocae059.
\bibitem{ref46}\label{r46} Styler WF 4th, Bethard S, Finan S, et al. Temporal Annotation in the Clinical Domain. Trans Assoc Comput Linguist 2014;2:143--154. doi:10.1162/tacl\_a\_00172.
\bibitem{ref47}\label{r47} Chu Z, Chen J, Chen Q, et al. TimeBench: A Comprehensive Evaluation of Temporal Reasoning Abilities in Large Language Models. Proc. 62nd Annu Meet Assoc Comput Linguist (ACL) 2024:1204--1228. doi:10.18653/v1/2024.acl-long.66.
\bibitem{ref48}\label{r48} Zolensky A, Jang KJ, Sabin J, et al. Speaker Role Identification in Clinical Conversations. Pac Symp Biocomput 2026;31:144--157. doi:10.1142/9789819824755\_0011.
\bibitem{ref49}\label{r49} Uzuner O, Bodnari A, Shen S, et al. Evaluating the state of the art in coreference resolution for electronic medical records. J Am Med Inform Assoc 2012;19(5):786--791. doi:10.1136/amiajnl-2011-000784.
\bibitem{ref50}\label{r50} Dechter R, Meiri I, Pearl J. Temporal constraint networks. Artif Intell 1991;49(1--3):61--95. doi:10.1016/0004-3702(91)90006-6.
\bibitem{ref51}\label{r51} Tsamardinos I, Vidal T, Pollack ME. CTP: A New Constraint-Based Formalism for Conditional, Temporal Planning. Constraints 2003;8(4):365--388. doi:10.1023/A:1025894003623.
\bibitem{ref52}\label{r52} Shahar Y. A framework for knowledge-based temporal abstraction. Artif Intell 1997;90(1--2):79--133. doi:10.1016/S0004-3702(96)00025-2.
\bibitem{ref53}\label{r53} Shahar Y, Miksch S, Johnson P. The Asgaard project: a task-specific framework for the application and critiquing of time-oriented clinical guidelines. Artif Intell Med 1998;14(1--2):29--51. doi:10.1016/S0933-3657(98)00015-3.
\bibitem{ref54}\label{r54} Sutton DR, Fox J. The syntax and semantics of the PROforma guideline modeling language. J Am Med Inform Assoc 2003;10(5):433--443. doi:10.1197/jamia.M1264.
\bibitem{ref55}\label{r55} Boxwala AA, Peleg M, Tu S, et al. GLIF3: a representation format for sharable computer-interpretable clinical practice guidelines. J Biomed Inform 2004;37(3):147--161. doi:10.1016/j.jbi.2004.04.002.
\bibitem{ref56}\label{r56} Nori H, King N, McKinney SM, et al. Capabilities of GPT-4 on Medical Challenge Problems. arXiv preprint 2023:2303.13375.
\bibitem{ref57}\label{r57} Builtjes L, Bosma J, Prokop M, et al. Leveraging open-source large language models for clinical information extraction in resource-constrained settings. JAMIA Open 2025;8(5):ooaf109. doi:10.1093/jamiaopen/ooaf109.
\bibitem{ref58}\label{r58} Ji Z, Lee N, Frieske R, et al. Survey of Hallucination in Natural Language Generation. ACM Comput Surv 2023;55(12):248. doi:10.1145/3571730.
\bibitem{ref59}\label{r59} Garcez A, Lamb LC. Neurosymbolic AI: the 3rd wave. Artif Intell Rev 2023;56:12387--12406. doi:10.1007/s10462-023-10448-w.
\bibitem{ref60}\label{r60} Prenosil GA, Weitzel TK, Bello SC, et al. Neuro-symbolic AI for auditable cognitive information extraction from medical reports. Commun Med 2025;5(1):491. doi:10.1038/s43856-025-01194-x.
\bibitem{ref61}\label{r61} Geng S, Josifoski M, Peyrard M, et al. Grammar-Constrained Decoding for Structured NLP Tasks without Finetuning. Proc. 2023 Conf Empir Methods Nat Lang Process (EMNLP) 2023:10932--10952. doi:10.18653/v1/2023.emnlp-main.674.
\bibitem{ref62}\label{r62} Geifman Y, El-Yaniv R. Selective Classification for Deep Neural Networks. Adv Neural Inf Process Syst (NeurIPS) 2017;30:4878--4887. papers.neurips.cc . https://papers.neurips.cc/paper/2017/hash/4a8423d5e91fda00bb7e46540e2b0cf1-Abstract.html
\bibitem{ref63}\label{r63} Guo C, Pleiss G, Sun Y, et al. On Calibration of Modern Neural Networks. Proc. 34th Int Conf Mach Learn (ICML) , PMLR 2017;70:1321--1330. proceedings.mlr.press/v70/guo17a . https://proceedings.mlr.press/v70/guo17a.html
\bibitem{ref64}\label{r64} Angelopoulos AN, Bates S. Conformal Prediction: A Gentle Introduction. Found Trends Mach Learn 2023;16(4):494--591. doi:10.1561/2200000101.
\bibitem{ref65}\label{r65} Angelopoulos AN, Bates S, Fisch A, et al. Conformal Risk Control. Proc. Int Conf Learn Represent (ICLR) 2024. arXiv:2208.02814.
\bibitem{ref66}\label{r66} Goddard K, Roudsari A, Wyatt JC. Automation bias: a systematic review of frequency, effect mediators, and mitigators. J Am Med Inform Assoc 2012;19(1):121--127. doi:10.1136/amiajnl-2011-000089.
\bibitem{ref67}\label{r67} US Food and Drug Administration (CDRH). Clinical Decision Support Software: Guidance for Industry and Food and Drug Administration Staff. Silver Spring (MD): FDA; revised January 2026. fda.gov (accessed 4 Sep 2026). https://www.fda.gov/regulatory-information/search-fda-guidance-documents/clinical-decision-support-software
\bibitem{ref68}\label{r68} Regulation (EU) 2024/1689 of the European Parliament and of the Council laying down harmonised rules on artificial intelligence (Artificial Intelligence Act). Off J Eur Union 2024;L 2024/1689. data.europa.eu/eli/reg/2024/1689/oj . http://data.europa.eu/eli/reg/2024/1689/oj
\bibitem{ref69}\label{r69} Apartsin A, Aperstein Y. Clinical Intent Extraction: A FHIR-Aligned Representation and the CIRCA Benchmark. Preprint; 2026. doi:10.13140/RG.2.2.28773.97760. CIRCA dataset, Zenodo: doi:10.5281/zenodo.22251272.
\bibitem{ref70}\label{r70} Laufer M, Aperstein Y, Apartsin A. Reliable Follow-Up Action and Date Extraction from Clinical Notes: A Hybrid Neural-Symbolic Approach. arXiv preprint; 2026. arXiv:2605.26560.
\bibitem{ref71}\label{r71} Van Tiem J, Cramer E, Iverson C, et al. Listening to the note: clinician perspectives on ambient artificial intelligence scribes in medical documentation. J Am Med Inform Assoc 2026;33(2):255--262. doi:10.1093/jamia/ocaf214.
\bibitem{ref72}\label{r72} Artsi Y, Sorin V, Glicksberg BS, et al. Large language models in real-world clinical workflows: a systematic review of applications and implementation. Front Digit Health 2025;7:1659134. doi:10.3389/fdgth.2025.1659134.
\end{thebibliography}